\input format.tex

\usepackage{graphicx}
\graphicspath{{cores/}}

\begin{document}

\vspace*{8mm}
\parindent0pt
%% 各章节
\setlength{\arrayrulewidth}{1pt}
\fontsize{9.3pt}{11pt}\selectfont
\color{gray2}

{\bf\sanhao 检测结果总览}

\vspace*{6mm}

\begin{LRaside}{肠道菌群概况}
\noindent\centering
\includegraphics[width=\linewidth]{gaikuang_orange.pdf}
\jiuhao
\noindent\begin{tabular}{@{}m{0.68cm}<{\centering}@{}m{0.68cm}<{\centering}@{}m{0.68cm}<{\centering}@{}m{0.68cm}<{\centering}@{}}
\includegraphics[width=.41cm]{sun_orange.pdf} & \tikz\draw[gray,fill=gray](0,0) circle[radius=4.1pt]; & \tikz\draw[gray,fill=gray](0,0) circle[radius=4.1pt]; & \tikz\draw[gray,fill=gray](0,0) circle[radius=4.1pt]; \\[-4pt]
\color{gray2}当次 & \color{gray2} & \color{gray2} & \color{gray2}
\\
\end{tabular}
\asidebreak % 分栏
1、您本次肠道菌群共检测出细菌种类169种，肠道菌群多样性
高，说明您的肠道菌群组成十分丰富，菌群失调风险低。\par\bigskip
2、您肠道中的有益菌含量
高，优于93.75{\%}的参考人群。\par\bigskip
3、您肠道内的有害菌含量
高于参考人群，请您注意调理肠道菌群。
\end{LRaside}

\begin{LRaside}{肠道菌群含量}
\noindent
\includegraphics[width=\linewidth]{fenbu_red.pdf}

\jiuhao
\noindent\begin{tabular}{@{}m{0.68cm}<{\centering}@{}m{0.68cm}<{\centering}@{}m{0.68cm}<{\centering}@{}m{0.68cm}<{\centering}@{}}
\includegraphics[width=.41cm]{sun_red.pdf} & \tikz\draw[gray,fill=gray](0,0) circle[radius=4.1pt]; & \tikz\draw[gray,fill=gray](0,0) circle[radius=4.1pt]; & \tikz\draw[gray,fill=gray](0,0) circle[radius=4.1pt]; \\[-4pt]
\color{gray2}当次 & \color{gray2} & \color{gray2} & \color{gray2}
\\
\end{tabular}
\asidebreak % 分栏
1、您肠道内的
布劳特氏菌属、栖粪杆菌属、粪球菌属、瘤胃球菌属、颤螺菌属等16项指标优于参考人群，有益于肠道健康；但拟杆菌属、多尔氏菌属、链球菌属、普雷沃氏菌属、萨特氏菌属等23项指标异常，不利于肠道健康，需引起注意，建议持续监测肠道健康。\par\bigskip
2、您的肠道内
检测出7种致病菌，其中脆弱拟杆菌、胶胨罗斯氏菌、产气荚膜梭菌、病原性大肠埃希氏菌、副流感嗜血杆菌等含量高于参考人群，可能会损害您的肠道健康，建议您持续监测肠道健康。
\end{LRaside}


\begin{LRaside}{肠道菌群营养功能分析}
\noindent
\includegraphics[width=\linewidth]{gongneng_orange.pdf}

\jiuhao
\noindent\begin{tabular}{@{}m{0.68cm}<{\centering}@{}m{0.68cm}<{\centering}@{}m{0.68cm}<{\centering}@{}m{0.68cm}<{\centering}@{}}
\includegraphics[width=.41cm]{sun_orange.pdf} & \tikz\draw[gray,fill=gray](0,0) circle[radius=4.1pt]; & \tikz\draw[gray,fill=gray](0,0) circle[radius=4.1pt]; & \tikz\draw[gray,fill=gray](0,0) circle[radius=4.1pt]; \\[-4pt]
\color{gray2}当次 & \color{gray2} & \color{gray2} & \color{gray2}
\\
\end{tabular}
\asidebreak %分栏
您肠道内参与维生素B1、维生素B2、维生素B3、维生素B5、维生素B6等10种营养物质代谢的菌群含量优于参考人群，有益于人体健康；参与胆碱代谢的菌群含量异常，请关注您的异常指标，持续监测肠道健康。
\end{LRaside}

颜色指示 \includegraphics[scale=1]{greencircle.pdf} {\bf 绿色表示健康}   \includegraphics[scale=1]{orangecircle.pdf}  {\bf 橙色表示需要关注}  \includegraphics[scale=1]{redcircle.pdf} {\bf 红色表示有风险}

\end{document}
